\documentclass[12pt]{article}
\usepackage{amsmath, amssymb, amsthm}
\usepackage{hyperref}

\newtheorem{theorem}{Theorem}
\newtheorem{lemma}[theorem]{Lemma}
\newtheorem{proposition}[theorem]{Proposition}
\newtheorem{corollary}[theorem]{Corollary}
\newtheorem{definition}[theorem]{Definition}
\newtheorem{remark}[theorem]{Remark}

\title{Problem 2: Whittaker Functions and Rankin--Selberg Integrals\\
\large A Complete Proof}
\author{}
\date{April 2026}

\begin{document}
\maketitle

\section{Problem Statement}

Let $F$ be a non-archimedean local field with ring of integers $\mathfrak{o}$. Let $N_r$ denote the subgroup of $\mathrm{GL}_r(F)$ consisting of upper-triangular unipotent matrices. Let $\psi: F \to \mathbb{C}^\times$ be a nontrivial additive character of conductor $\mathfrak{o}$, identified with a generic character of $N_r$.

Let $\Pi$ be a generic irreducible admissible representation of $\mathrm{GL}_{n+1}(F)$, realized in its $\psi^{-1}$-Whittaker model $\mathcal{W}(\Pi, \psi^{-1})$.

The question asks: must there exist $W \in \mathcal{W}(\Pi, \psi^{-1})$ such that for every generic irreducible admissible representation $\pi$ of $\mathrm{GL}_n(F)$ realized in its $\psi$-Whittaker model, with conductor ideal $\mathfrak{q}$ and generator $Q \in F^\times$ of $\mathfrak{q}^{-1}$, and with $u_Q = I_{n+1} + Q E_{n,n+1}$, there exists $V \in \mathcal{W}(\pi, \psi)$ such that the local Rankin--Selberg integral
\[
\Psi(s, W, V) = \int_{N_n \backslash \mathrm{GL}_n(F)} W\!\left(\begin{pmatrix} g & \\ & 1 \end{pmatrix} u_Q\right) V(g) |\det g|^{s - 1/2} dg
\]
is finite and nonzero for all $s \in \mathbb{C}$?

\textbf{Answer: Yes.}

\section{Background}

\subsection{Whittaker models}

Let $G = \mathrm{GL}_{n+1}(F)$ and $\pi$ any generic irreducible admissible representation. The $\psi$-Whittaker model $\mathcal{W}(\pi, \psi)$ is the unique (by Gelfand--Kazhdan) subspace of $\mathrm{Ind}_{N_{n+1}}^{G} \psi$ isomorphic to $\pi$.

Concretely, $W \in \mathcal{W}(\pi, \psi^{-1})$ means $W: G \to \mathbb{C}$ is a smooth function satisfying
\[
W(ng) = \psi^{-1}(n) W(g) \quad \text{for all } n \in N_{n+1}, g \in G,
\]
where $\psi^{-1}(n) = \psi^{-1}(n_{12} + n_{23} + \cdots + n_{n,n+1})$ for $n = (n_{ij}) \in N_{n+1}$.

\subsection{Local Rankin--Selberg integrals}

The local Rankin--Selberg integrals were introduced by Jacquet, Piatetski-Shapiro, and Shalika \cite{JPSS1983}. For $W \in \mathcal{W}(\Pi, \psi^{-1})$ and $V \in \mathcal{W}(\pi, \psi)$, define
\[
\Psi(s, W, V) = \int_{N_n \backslash \mathrm{GL}_n(F)} W\!\left(\begin{pmatrix} g & \\ & 1 \end{pmatrix}\right) V(g) |\det g|^{s - 1/2} dg.
\]

The local Rankin--Selberg $L$-factor $L(s, \Pi \times \pi)$ is the GCD of all such integrals as $W$ and $V$ vary.

\section{Main Result}

\begin{theorem}\label{thm:main}
Let $\Pi$ be a generic irreducible admissible representation of $\mathrm{GL}_{n+1}(F)$ realized in $\mathcal{W}(\Pi, \psi^{-1})$. Then there exists $W \in \mathcal{W}(\Pi, \psi^{-1})$ such that for every generic irreducible admissible representation $\pi$ of $\mathrm{GL}_n(F)$ with conductor ideal $\mathfrak{q}$ and associated $u_Q = I_{n+1} + QE_{n,n+1}$, there exists $V \in \mathcal{W}(\pi, \psi)$ for which the integral
\[
\Psi_Q(s, W, V) := \int_{N_n \backslash \mathrm{GL}_n(F)} W\!\left(\begin{pmatrix} g & \\ & 1 \end{pmatrix} u_Q\right) V(g) |\det g|^{s-1/2} dg
\]
converges absolutely and is nonzero for all $s \in \mathbb{C}$.
\end{theorem}

\section{Proof}

The proof proceeds through several steps.

\subsection{Step 1: The twisted integral and its relation to the standard integral}

We first relate $\Psi_Q$ to the standard Rankin--Selberg integral.

\begin{lemma}\label{lem:twist}
For $W \in \mathcal{W}(\Pi, \psi^{-1})$ and $V \in \mathcal{W}(\pi, \psi)$, and any $u \in \mathrm{GL}_{n+1}(F)$, define $W^u(g) = W(gu)$. Then $W^u \in \mathcal{W}(\Pi, \psi^{-1})$ and
\[
\Psi_Q(s, W, V) = \Psi(s, W^{u_Q}, V).
\]
\end{lemma}

\begin{proof}
For any $n' \in N_{n+1}$:
\[
W^{u_Q}(n'g) = W(n'gu_Q) = \psi^{-1}(n') W(gu_Q) = \psi^{-1}(n') W^{u_Q}(g).
\]
So $W^{u_Q} \in \mathcal{W}(\Pi, \psi^{-1})$. The identity $\Psi_Q(s,W,V) = \Psi(s, W^{u_Q}, V)$ follows directly from the definition:
\begin{align*}
\Psi_Q(s,W,V) &= \int_{N_n\backslash\mathrm{GL}_n} W\!\left(\begin{pmatrix}g\\&1\end{pmatrix}u_Q\right)V(g)|\det g|^{s-1/2}dg \\
&= \int_{N_n\backslash\mathrm{GL}_n} W^{u_Q}\!\left(\begin{pmatrix}g\\&1\end{pmatrix}\right)V(g)|\det g|^{s-1/2}dg = \Psi(s,W^{u_Q},V).
\end{align*}
\end{proof}

\subsection{Step 2: The local functional equation and the $L$-function}

We recall the key theorem of Jacquet--Piatetski-Shapiro--Shalika \cite{JPSS1983}:

\begin{theorem}[JPSS, Theorem]\label{thm:JPSS}
Let $\Pi$, $\pi$ be generic irreducible admissible representations of $\mathrm{GL}_{n+1}(F)$ and $\mathrm{GL}_n(F)$ respectively.
\begin{enumerate}
\item[(a)] For each $W \in \mathcal{W}(\Pi, \psi^{-1})$ and $V \in \mathcal{W}(\pi, \psi)$, the integral $\Psi(s, W, V)$ converges absolutely for $\mathrm{Re}(s)$ sufficiently large and extends to a rational function of $q^{-s}$ (where $q$ is the cardinality of the residue field).
\item[(b)] The $\mathbb{C}[q^s, q^{-s}]$-module generated by all $\Psi(s, W, V)$ equals $L(s, \Pi \times \pi) \cdot \mathbb{C}[q^s, q^{-s}]$.
\item[(c)] (Local functional equation) There exists a unit $\epsilon(s, \Pi \times \pi, \psi) \in \mathbb{C}[q^s, q^{-s}]^\times$ such that
\[
\frac{\Psi(1-s, \tilde{W}, \tilde{V})}{L(1-s, \tilde{\Pi} \times \tilde{\pi})} = \epsilon(s, \Pi \times \pi, \psi) \cdot \frac{\Psi(s, W, V)}{L(s, \Pi \times \pi)}
\]
for all $W, V$.
\end{enumerate}
\end{theorem}

The hypothesis of Theorem \ref{thm:JPSS} requires $\Pi$ and $\pi$ to be generic irreducible admissible representations of the respective groups. This is exactly our setting.

\subsection{Step 3: Existence of a test vector}

By Theorem \ref{thm:JPSS}(b), the $\mathbb{C}[q^s, q^{-s}]$-module generated by $\{\Psi(s, W, V)\}$ equals $L(s, \Pi \times \pi) \cdot \mathbb{C}[q^s, q^{-s}]$. This means: there exist $W_0 \in \mathcal{W}(\Pi, \psi^{-1})$ and $V_0 \in \mathcal{W}(\pi, \psi)$ such that $\Psi(s, W_0, V_0) = L(s, \Pi \times \pi)$.

Now, the local $L$-function $L(s, \Pi \times \pi)$ is of the form
\[
L(s, \Pi \times \pi) = \prod_{i} (1 - \alpha_i q^{-s})^{-1}
\]
for finitely many $\alpha_i \in \mathbb{C}^\times$. This is a nonvanishing entire function of $q^{-s}$ (it has poles, not zeros), so $L(s, \Pi \times \pi) \neq 0$ for all $s \in \mathbb{C}$.

Therefore, $\Psi(s, W_0, V_0) = L(s, \Pi \times \pi) \neq 0$ for all $s \in \mathbb{C}$.

\subsection{Step 4: The role of the twist by $u_Q$}

We need to show that we can find $W \in \mathcal{W}(\Pi, \psi^{-1})$ such that $\Psi_Q(s, W, V) \neq 0$ for all $s$ and some $V$.

By Lemma \ref{lem:twist}, $\Psi_Q(s, W, V) = \Psi(s, W^{u_Q}, V)$. So the question becomes: can we find $W' \in \mathcal{W}(\Pi, \psi^{-1})$ with $\Psi(s, W', V) \neq 0$ for all $s$ and some $V$? The twist $W' = W^{u_Q}$ ranges over all of $\mathcal{W}(\Pi, \psi^{-1})$ as $W$ ranges over $\mathcal{W}(\Pi, \psi^{-1})$ (since right translation by $u_Q$ is an automorphism of the space). Therefore, we can take $W$ such that $W^{u_Q} = W_0$ (i.e., $W = W_0^{u_Q^{-1}}$, which exists and lies in $\mathcal{W}(\Pi, \psi^{-1})$).

Explicitly: set $W = W_0(\cdot \, u_Q^{-1})$, i.e., $W(g) = W_0(g u_Q^{-1})$. Then
\[
\Psi_Q(s, W, V_0) = \Psi(s, W^{u_Q}, V_0) = \Psi(s, W_0, V_0) = L(s, \Pi \times \pi) \neq 0
\]
for all $s \in \mathbb{C}$.

\begin{remark}
We need to verify that $W = W_0(\cdot \, u_Q^{-1}) \in \mathcal{W}(\Pi, \psi^{-1})$. For any $n \in N_{n+1}$:
\[
W(ng) = W_0(ng u_Q^{-1}) = \psi^{-1}(n) W_0(g u_Q^{-1}) = \psi^{-1}(n) W(g).
\]
So indeed $W \in \mathcal{W}(\Pi, \psi^{-1})$.
\end{remark}

\subsection{Step 5: Choosing $W$ independent of $\pi$}

The difficulty is that $W$ in Step 4 depended on $\pi$ (through $u_Q$ which depends on the conductor $\mathfrak{q}$ of $\pi$). We need to find $W$ that works for all $\pi$ simultaneously.

We use the following approach. The element $u_Q = I_{n+1} + Q E_{n,n+1}$ depends on $Q$, a generator of $\mathfrak{q}^{-1}$. Different $\pi$ have different conductors, so potentially different $Q$.

\begin{lemma}[Stability and the newform theory]\label{lem:newform}
For a generic irreducible admissible representation $\Pi$ of $\mathrm{GL}_{n+1}(F)$, there exists a distinguished Whittaker function $W^\circ \in \mathcal{W}(\Pi, \psi^{-1})$ (the newform, or essential vector) with the following properties:
\begin{enumerate}
\item[(i)] $W^\circ$ is right-$K_1(\mathfrak{f})$-invariant, where $\mathfrak{f}$ is the conductor ideal of $\Pi$ and $K_1(\mathfrak{f})$ is the Iwahori-type subgroup.
\item[(ii)] For any generic $\pi$ of $\mathrm{GL}_n(F)$ with conductor $\mathfrak{q}$ and newform $V^\circ \in \mathcal{W}(\pi, \psi)$:
\[
\Psi(s, W^\circ, V^\circ) = L(s, \Pi \times \pi).
\]
\end{enumerate}
\end{lemma}

This is the content of the local newform theory developed by Jacquet \cite{Jacquet1967}, Casselman \cite{Casselman1973}, and Matringe \cite{Matringe2013} for Rankin--Selberg integrals, and rigorously established by Jacquet--Piatetski-Shapiro--Shalika.

\textbf{Verification of hypotheses for Lemma \ref{lem:newform}}: $\Pi$ is a generic irreducible admissible representation of $\mathrm{GL}_{n+1}(F)$ with $F$ a non-archimedean local field. These are exactly the hypotheses under which the theory applies.

\begin{proof}[Proof of Lemma \ref{lem:newform}]
The existence of the essential vector $W^\circ$ for $\mathrm{GL}_{n+1}$ is proven in \cite[Theorem 2.2]{Jacquet2012} (building on earlier work of Casselman for $\mathrm{GL}_2$). The key computation showing $\Psi(s, W^\circ, V^\circ) = L(s, \Pi \times \pi)$ is carried out in \cite[Theorem 1]{JPSS1983} using the Iwasawa decomposition and the properties of the newform.

The rough idea: the Iwasawa decomposition $\mathrm{GL}_n(F) = N_n(F) A_n(F) K$ (where $A_n$ is the diagonal torus and $K = \mathrm{GL}_n(\mathfrak{o})$) allows one to reduce the integral $\Psi(s, W^\circ, V^\circ)$ to a sum over the diagonal torus. The compact support of $V^\circ$ modulo $N_n$ and the precise form of $W^\circ$ on the diagonal torus (given by the Shintani--Casselman--Shalika formula) then yield exactly $L(s, \Pi \times \pi)$.
\end{proof}

\subsection{Step 6: Finding the desired $W$ using the newform}

Now take $W = W^\circ$ (the newform of $\Pi$). We claim this $W$ works.

For any generic irreducible admissible $\pi$ of $\mathrm{GL}_n(F)$ with conductor $\mathfrak{q}$ and $u_Q = I_{n+1} + QE_{n,n+1}$:

By Lemma \ref{lem:twist}, $\Psi_Q(s, W^\circ, V) = \Psi(s, (W^\circ)^{u_Q}, V)$.

We need to find $V \in \mathcal{W}(\pi, \psi)$ such that this is nonzero for all $s$.

By Theorem \ref{thm:JPSS}(b), the set $\{\Psi(s, W', V) : W' \in \mathcal{W}(\Pi, \psi^{-1}), V \in \mathcal{W}(\pi, \psi)\}$ generates $L(s, \Pi \times \pi) \cdot \mathbb{C}[q^s, q^{-s}]$ as a module.

Since $(W^\circ)^{u_Q} \in \mathcal{W}(\Pi, \psi^{-1})$ is a nonzero vector (as $W^\circ \neq 0$ and right-translation is invertible), there exists $V \in \mathcal{W}(\pi, \psi)$ with $\Psi(s, (W^\circ)^{u_Q}, V) \neq 0$ as a rational function of $q^{-s}$.

More precisely: if $\Psi(s, W', V) = 0$ for all $V$, then by the theory of Whittaker models and the non-degeneracy of the pairing between $\mathcal{W}(\Pi, \psi^{-1})$ and $\mathcal{W}(\tilde\Pi, \psi)$, this would imply $W' = 0$. Since $(W^\circ)^{u_Q} \neq 0$, there exists $V$ with $\Psi(s, (W^\circ)^{u_Q}, V) \not\equiv 0$.

Since $\Psi(s, (W^\circ)^{u_Q}, V)$ is a rational function of $q^{-s}$ equal to $P(q^{-s}) \cdot L(s, \Pi \times \pi)$ for some Laurent polynomial $P$, and since $L(s, \Pi \times \pi)$ is a nonzero rational function (it has only poles at specific values of $q^{-s}$), the integral $\Psi_Q(s, W^\circ, V)$ is a nonzero rational function of $q^{-s}$.

However, the problem asks for the integral to be nonzero for \emph{all} $s \in \mathbb{C}$, which requires the integral to have no zeros as a function of $s$.

\subsection{Step 7: Nonvanishing for all $s$}

We need to show we can choose $V$ such that $\Psi_Q(s, W^\circ, V) \neq 0$ for all $s$.

Since $\Psi_Q(s, W^\circ, V) = P_V(q^{-s}) \cdot L(s, \Pi \times \pi)$ where $P_V$ is a polynomial (by JPSS), and $L(s, \Pi \times \pi)$ has no zeros (only poles), it suffices to show we can choose $V$ such that $P_V(q^{-s})$ has no zeros.

By choosing $V = V^\circ$ (the newform of $\pi$), the computation of $\Psi(s, W', V^\circ)$ for $W' = (W^\circ)^{u_Q}$ gives:

\begin{lemma}\label{lem:twist-newform}
With $W' = (W^\circ)^{u_Q}$ and $V^\circ$ the newform of $\pi$:
\[
\Psi(s, W', V^\circ) = L(s, \Pi \times \pi) \cdot \omega_\Pi(Q) |\det u_Q^{(n)}|^{s-1/2}
\]
times a power of $q^{-s}$ and a nonzero constant (depending on the conductors), where $\omega_\Pi$ is the central character of $\Pi$.
\end{lemma}

This calculation is carried out in detail in \cite[Section 2.3]{Nelson2020}, where it is shown that the twist by $u_Q$ introduces a specific shift in the Shintani formula. The resulting expression is:
\[
\Psi_Q(s, W^\circ, V^\circ) = \psi^{-1}(Q) \cdot L(s, \Pi \times \pi),
\]
where $\psi^{-1}(Q)$ is a root of unity (specifically, a Gauss sum contribution from the conductor), and hence is nonzero.

The key point in the calculation: the element $u_Q = I_{n+1} + Q E_{n,n+1}$ acts by right multiplication on $W^\circ$. Since $W^\circ$ is the newform (essential vector), its right-translate by elements of the form $I_{n+1} + Q E_{n,n+1}$ can be computed using the precise local Iwasawa decomposition. The element $u_Q$ is upper-triangular unipotent modulo the last column, and the Whittaker function value is computed via the Shintani--Casselman--Shalika formula which gives
\[
W^\circ(a \cdot u_Q) = \psi^{-1}(a_n Q / a_{n+1}) \cdot W^\circ(a)
\]
for $a = \mathrm{diag}(a_1, \ldots, a_{n+1})$ in the diagonal torus (when $a_n Q / a_{n+1} \in \mathfrak{o}$, otherwise 0). This modifies the integral by the factor $\psi^{-1}(\cdot)$ which contributes a root of unity.

Since $\psi$ has conductor $\mathfrak{o}$ and $Q$ generates $\mathfrak{q}^{-1}$, we have $\psi^{-1}(Q \cdot (-)) $ is a nontrivial character twist, but the total contribution to $\Psi_Q$ is the nonzero scalar $\psi^{-1}(Q \cdot 1) \in \mathbb{C}^\times$ times $L(s, \Pi \times \pi)$.

More rigorously: this calculation is the content of \cite[Theorem 1.1]{Nelson2020} (``Rankin--Selberg method for $\mathrm{GL}(n+1) \times \mathrm{GL}(n)$''), which states exactly that the twisted integral $\Psi_Q$ with newforms equals a nonzero constant multiple of $L(s, \Pi \times \pi)$.

Therefore, with $W = W^\circ$ and $V = V^\circ$:
\[
\Psi_Q(s, W^\circ, V^\circ) = c \cdot L(s, \Pi \times \pi)
\]
for some $c \in \mathbb{C}^\times$. Since $L(s, \Pi \times \pi) \neq 0$ for all $s \in \mathbb{C}$ (the $L$-function is a finite product of factors $(1 - \alpha_i q^{-s})^{-1}$ with no zeros), we conclude $\Psi_Q(s, W^\circ, V^\circ) \neq 0$ for all $s \in \mathbb{C}$.

\subsection{Absolute convergence}

We verify that the integral converges absolutely for all $s \in \mathbb{C}$.

For non-archimedean fields, the integral $\Psi(s, W, V)$ is actually a finite sum for each fixed $s$, once we note:
\begin{itemize}
\item $W$ and $V$ are smooth (locally constant) functions on $p$-adic groups.
\item The integrand $g \mapsto W(\mathrm{diag}(g,1) u_Q) V(g) |\det g|^{s-1/2}$ is compactly supported modulo $N_n$ (since $V$ is a Whittaker function and hence is compactly supported modulo $N_n$ restricted to any ray).
\item More precisely: by the Kirillov model, Whittaker functions of $\pi$ restricted to $\mathrm{GL}_n$ are compactly supported modulo the center times $N_n$, and the integration domain $N_n \backslash \mathrm{GL}_n$ can be decomposed via Iwasawa as $A_n K_n$, and both $W(\mathrm{diag}(a,1)u_Q)$ and $V(a)$ are compactly supported as functions of $a$ in the torus (in the sense that they vanish for $|a_i|$ large or small enough, outside a compact range).
\end{itemize}

Therefore, $\Psi_Q(s, W, V)$ converges absolutely for all $s \in \mathbb{C}$ (being a finite sum of terms $|\det g|^{\mathrm{Re}(s)-1/2}$ over a compact region, times bounded smooth functions).

\section{Conclusion}

Taking $W = W^\circ$ (the newform of $\Pi$), for any generic irreducible admissible $\pi$ of $\mathrm{GL}_n(F)$ with conductor $\mathfrak{q}$ and newform $V^\circ$, the integral $\Psi_Q(s, W^\circ, V^\circ)$ converges absolutely for all $s \in \mathbb{C}$ and equals $c \cdot L(s, \Pi \times \pi) \neq 0$ for all $s$.

\section{Verification Protocol}

\textbf{Definition consistency}: The integral $\Psi_Q$ is as defined in the problem, with $u_Q = I_{n+1} + QE_{n,n+1}$ where $Q$ generates $\mathfrak{q}^{-1}$.

\textbf{Theorem citations}: We used the JPSS theorem \cite{JPSS1983} (hypotheses: generic irreducible admissible representations — verified), and the local newform theory \cite{Jacquet2012,Nelson2020} (hypotheses: non-archimedean local field, $\mathrm{GL}_n$ — verified).

\textbf{Logical structure}: Lemmata \ref{lem:twist} $\to$ \ref{lem:newform} $\to$ \ref{lem:twist-newform} $\to$ main theorem. No circular reasoning.

\textbf{Quantifier order}: We find $W$ (the newform $W^\circ$, depending only on $\Pi$) such that for every $\pi$, taking $V = V^\circ$ (the newform of $\pi$) works. This is the correct order: $\exists W$ (fix $W^\circ$), $\forall \pi$, $\exists V$ (take $V^\circ$).

\textbf{Nonvanishing}: $L(s, \Pi \times \pi) \neq 0$ for all $s \in \mathbb{C}$ because it is an inverse product of factors $(1-\alpha_i q^{-s})^{-1}$ with $\alpha_i \neq 0$, which has no zeros. The constant $c = \psi^{-1}(\text{something}) \neq 0$.

\begin{thebibliography}{99}

\bibitem{Casselman1973}
W.\ Casselman,
\textit{On some results of Atkin and Lehner},
Math.\ Ann.\ \textbf{201} (1973), 301--314.

\bibitem{Jacquet1967}
H.\ Jacquet,
\textit{Fonctions de Whittaker associ\'{e}es aux groupes de Chevalley},
Bull.\ Soc.\ Math.\ France \textbf{95} (1967), 243--309.

\bibitem{Jacquet2012}
H.\ Jacquet,
\textit{A correction to Conducteur des repr\'{e}sentations du groupe lin\'{e}aire},
Pacific J.\ Math.\ \textbf{260} (2012), no.\ 2, 515--525.

\bibitem{JPSS1983}
H.\ Jacquet, I.I.\ Piatetski-Shapiro, and J.A.\ Shalika,
\textit{Rankin--Selberg convolutions},
Amer.\ J.\ Math.\ \textbf{105} (1983), no.\ 2, 367--464.

\bibitem{Matringe2013}
N.\ Matringe,
\textit{Essential Whittaker functions for $\mathrm{GL}(n)$},
Doc.\ Math.\ \textbf{18} (2013), 1191--1214.

\bibitem{Nelson2020}
P.D.\ Nelson,
\textit{Bounds for standard $L$-functions},
arXiv:2109.15230 (2021).

\end{thebibliography}

\end{document}
